% Generated draft: requires mathematical and source review.
\documentclass[11pt]{article}
\usepackage[margin=25mm]{geometry}
\usepackage{fontspec}
\setmainfont{Latin Modern Roman}
\usepackage{amsmath,amssymb,mathtools,unicode-math}
\setmathfont{Latin Modern Math}
\usepackage{enumitem,needspace,xurl,hyperref}
\setlength{\parindent}{0pt}
\setlength{\parskip}{5pt}
\setlength{\emergencystretch}{3em}
\newcommand{\N}{\mathbb{N}}
\newcommand{\Z}{\mathbb{Z}}
\newcommand{\Q}{\mathbb{Q}}
\newcommand{\R}{\mathbb{R}}
\newcommand{\C}{\mathbb{C}}
\newcommand{\F}{\mathbb{F}}
\newcommand{\A}{\mathbb{A}}
\newcommand{\PP}{\mathbb{P}}
\newcommand{\E}{\mathbb{E}}
\newcommand{\eps}{\varepsilon}
\newcommand{\dd}{\,\mathrm{d}}
\begin{document}
\section*{UnsolvedMath ID: 30004749}
\textbf{Source problem code:} catalogue entry
\subsection{Definitions and mathematical statement}
Let \(K\) be a number field and let \(X\) be a smooth projective curve over \(K\), as specified in the supplied statement. For a closed point \(P\in X\), let \(\kappa(P)\) denote its residue field and define its degree over \(K\) by
\[
\deg_K(P)=[\kappa(P):K].
\]
For an integer \(d\), put
\[
S_d(X/K)=\{P\in X:\ P\text{ is a closed point and }[\kappa(P):K]=d\}.
\]
Under this interpretation, \(S_d(X/K)=\varnothing\) for \(d\leq 0\).

\textbf{Problem.} For each number field \(K\), smooth projective curve \(X/K\), and positive integer \(d\), determine whether \(S_d(X/K)\) is finite or infinite; that is, characterize the inputs \((K,X,d)\) for which infinitely many closed points have degree exactly \(d\).

\textbf{Unresolved qualifications.} The source does not explicitly define a degree-\(d\) point; the formulation above uses closed points of exact residue-field degree, rather than effective divisors of degree \(d\) or points of degree at most \(d\). The original wording uses ``nice curve,'' whereas the supplied cleaned statement specifies smoothness and projectivity without explicitly requiring geometric integrality or connectedness. No such additional hypothesis or genus restriction is imposed here. The source asks for a finiteness classification but supplies no proposed criterion or precise required form of that classification.
\subsection{Short English statement}
For a smooth projective curve over a number field and a positive integer \(d\), determine when there are finitely or infinitely many closed points whose residue fields have degree exactly \(d\). The meaning of ``degree-\(d\) point'' and the hypotheses implicit in the original phrase ``nice curve'' require confirmation.
\subsection{Sources}
\begin{itemize}
\item Supplied local UnsolvedMath registry record, permanent ID \(30004749\), title ``Degree-\(d\) Points on Algebraic Curves,'' exact source problem code \texttt{catalogue entry}. The record gives proposed year 2021 and identifies no individual proposer.
\item Supplied source citation: \emph{Explicit Methods in Number Theory (2022)}, Oberwolfach Reports, pp.\ 1805--1858, DOI \texttt{10.4171/owr/2021/34}. \url{https://doi.org/10.4171/owr/2021/34}. The supplied background describes these reports as documenting mathematical workshops and their open-problem discussions.
\item Supplied literature source, labelled ``Kadets--Vogt'': \url{https://arxiv.org/abs/2208.01067}.
\item Supplied literature source, labelled ``Oberwolfach report'': \url{https://ems.press/content/serial-article-files/46913}.
\item Literature assessment dated exactly \texttt{2026-08-22}: the supplied curation labels the problem open and reports that important low-degree configurations have been classified, especially in degrees two and three and selected higher degrees, while the stated extension of Mordell remains unresolved. This is supplied dated curation derived from the OWR record, not an independently verified assessment; time-sensitive status and cited claims require independent verification.
\item The registry labels the statement \texttt{corrected\_verified} and reports a manual semantic, typographic, and LaTeX review of the cleaned statement. This is supplied review metadata, not an independent verification of the source or mathematical formulation.
\end{itemize}
\end{document}
